%% sections/01-introduction.tex

\section{Introduction}
\label{sec:introduction}

Campaign design for tabletop role-playing games (TRPGs) is one of the least
systematized domains of game design practice. Adventure design has received
meaningful treatment in both practitioner literature and academic
research~\cite{tychsen2006role, dormans2010}, and encounter design is reasonably
well understood through the lens of challenge calibration and
pacing~\cite{dungeons2014players}. But the \textit{campaign} --- the multi-session arc
that structures a series of adventures into a coherent narrative whole --- lacks an
equivalent body of validated design knowledge.

The consequences are visible in practice. Practitioner surveys suggest that most
long-form TRPG campaigns fail to reach a designed conclusion~\cite{perkins2017dungeon}.
Sessions diverge from the arc's design intent; central questions go unanswered;
NPCs designed for the finale are never reached; the campaign ends when players
become bored rather than when the narrative is complete. These failures are typically
attributed to scheduling difficulties, player attrition, or DM burnout. We propose
a structural cause that has received less attention: campaigns without explicit
completion architecture cannot self-terminate. When there is no designed state
that constitutes ``complete'', the campaign continues until external forces stop
it rather than until its narrative logic is exhausted.

This paper presents the \textit{campaign spine} --- a five-component operational
document for multi-session TRPG campaign design --- and reports its validation
across 7 complete Dungeons \& Dragons 5e campaigns of 7 sessions each, run in the
Marathon AI-simulated TRPG workshop. The spine's core claim is that reliable
campaign completion requires explicit design of five components:
(a)~a \textit{central question} answerable by party behavior;
(b)~a \textit{4-hint delivery schedule} that builds the party's capacity to
recognize and reach the optimal ending;
(c)~\textit{branching end-conditions} with explicit route requirements;
(d)~a \textit{PC-spotlight schedule} that distributes character arc moments
across the campaign; and
(e)~a \textit{PC-authored finale deliverable specification} that accumulates
across sessions and is produced verbatim at the finale.

\subsection{The Completion Problem in Long-Form TRPG}

The problem of campaign arc completion is well-known to practitioners. A TRPG
campaign that begins with a clear premise and engaged players still faces structural
hazards across its arc: sessions that produce divergent narrative threads; players
who invest in sub-goals the campaign spine cannot accommodate; DM attention that
concentrates on session-level design at the expense of arc-level coherence;
and endings that are improvised rather than designed, producing conclusions that
feel abrupt, arbitrary, or unearned.

Published campaign modules address this partly through adventure sequence (each
adventure leads to the next) but rarely through explicit arc completion criteria.
A published campaign specifying that ``the PCs must confront the villain in
adventure 7'' does not guarantee that the confrontation produces emotional
completeness: it specifies a plot endpoint, not an emotional arc endpoint.
The distinction matters. A plot endpoint is satisfied when the villain is
defeated. An emotional arc endpoint requires that the party's relationship
to the campaign's central question has been answered by their own actions ---
that they have earned the ending, not merely arrived at it.

The AI-simulated TRPG context makes this structural requirement especially
visible. Without human social dynamics, without DM affect-reading and
real-time improvisation, without the conversational momentum that human
players generate, an AI-simulated campaign has no mechanism for producing
emotional completeness that is not built into the design. If the design does
not specify when the campaign is complete and what that completion requires,
no amount of AI capability can supply it. The campaign spine emerged precisely
from this constraint: it is the minimum structural specification that allows
an AI-simulated campaign to reach genuine emotional completion.

\subsection{The Campaign Spine: A Structural Solution}

The campaign spine addresses the completion problem through five design
components that operate at different timescales within the campaign.
The central question provides the arc's organizing principle across all
7 sessions. The hint schedule creates a convergence trajectory: the party
receives pieces of insight at designed intervals, each piece building toward
the finale's optimal condition. The branching end-conditions ensure the
campaign can terminate gracefully at any engagement level rather than
requiring either the hardest route or nothing. The spotlight schedule
distributes PC arc moments so that every character's story has a dedicated
session rather than competing for space in the finale. The deliverable
specification creates a thread of continuity from session to session:
a document, speech, or act that the party accumulates material for across
the campaign and produces in full at the end.

These five components are not independent. They interact in specific ways
that the design must account for, and their interactions are as important
as the components themselves. Section~\ref{sec:components} details both
the components and their designed interactions.

\subsection{Evidence Summary}

Across 7 complete campaigns (49 sessions total) in the Marathon workshop:

\begin{itemize}
  \item \textbf{7/7 campaigns completed} within 7 sessions (100\% completion rate).
  \item \textbf{5/7 campaigns achieved Route D or E}, the end-conditions
    requiring the highest sustained party engagement with the spine's design.
  \item \textbf{28/28 hints delivered} across all campaigns (100\% hint
    delivery rate). Three campaigns used documented recovery paths; no hint was
    permanently undelivered.
  \item \textbf{7/7 PC-authored deliverables produced} verbatim in finale sessions.
  \item \textbf{Mean sessions-to-first-binding-PASS: 4.5} across all campaigns
    (range: 4--6).
  \item \textbf{Archetype independence confirmed}: equivalent outcomes across grief,
    covenant, faith/authority, siege-resource, professional-code, duty-spine, and
    no-designed-stakes campaign types.
\end{itemize}

\subsection{Contributions}

This paper makes four primary contributions:

\begin{enumerate}
  \item \textbf{The spine specification:} Five components with implementation
    guidance, illustrated with examples from the corpus.

  \item \textbf{Cross-campaign validation:} Outcome data for 7 complete
    campaigns across diverse archetypes, with full route, hint-delivery,
    and deliverable evidence.

  \item \textbf{Component interaction analysis:} Evidence on which components
    most contribute to finale reliability and how components depend on one another.

  \item \textbf{Archetype independence:} Evidence that the spine's structural
    features generalize across fundamentally different motivational architectures,
    including parties with no designed personal stakes.
\end{enumerate}

\subsection{Paper Organization}

Section~\ref{sec:related} situates the campaign spine in the literature on
narrative structure, quest design, and interactive drama. Section~\ref{sec:components}
presents the five components, their design rationale, and their interactions.
Section~\ref{sec:evaluation} reports the cross-campaign quantitative evidence,
including campaign outcome table, hint delivery mapping, route distribution,
and sessions-to-first-binding-PASS analysis. Section~\ref{sec:discussion}
interprets the findings, addresses limitations, and describes future work.
Section~\ref{sec:conclusion} concludes.
